Với lịch 40 epoch, Swin-T scratch đạt accuracy test $60{,}00 \pm 0{,}47$\%, cao hơn rõ rệt so với 2,19\% của lịch 20 epoch. Đường cong train và validation đều cải thiện, cho thấy mô hình đã học được cách phân biệt các lớp. Tuy vậy, accuracy cuối lịch trên train gần 90\%, còn validation khoảng 60\%: khớp dữ liệu train vẫn dễ hơn phân loại ảnh mới. Các checkpoint scratch tốt nhất được chọn trong khoảng epoch 36--37.

Swin-T pretrained đạt $86{,}45 \pm 0{,}19$\% dưới cùng lịch. Để so sánh, lấy accuracy của bản pretrained trừ bản scratch ở từng seed, rồi tính trung bình ba chênh lệch. Kết quả là $26{,}45 \pm 0{,}40$ điểm phần trăm; bản pretrained cao hơn ở cả ba cặp. Như vậy, ngay khi bản scratch đã học được, khởi tạo từ ImageNet vẫn đem lại lợi ích lớn. Hai bản có cùng kiến trúc, cập nhật toàn bộ tham số và dùng cùng lịch, nên phép so sánh này rõ hơn về vai trò của khởi tạo so với lịch 20 epoch.

\begin{table}[htbp]\centering\small
\setlength{\tabcolsep}{4pt}
\begin{tabular}{llrrrr}
\toprule
Kiến trúc & Chế độ & Tổng (M) & Học (M) & Giây & MiB \\
\midrule
Swin-T & Scratch & 27{,}596 & 27{,}596 & 428{,}1 $\pm$ 1{,}6 & 2453{,}2 $\pm$ 3{,}7 \\
Swin-T & Full & 27{,}596 & 27{,}596 & 429{,}7 $\pm$ 0{,}8 & 2456{,}7 $\pm$ 2{,}2 \\
\bottomrule
\end{tabular}

\caption{Chi phí lịch 40 epoch; thời gian train và peak allocated memory là trung bình và SD mẫu qua ba seed.}
\end{table}

Chi phí hai bản gần nhau: thời gian train đều khoảng 429 giây và bộ nhớ cấp phát cực đại khoảng 2455 MiB. Khác biệt chính trong đối chứng này nằm ở chất lượng phân loại, không phải số tham số hoặc ngân sách tính toán.
